problem:  
Let \( G \) be a connected semisimple Lie group with finite center, \( H\subset G \) a maximal compact subgroup, and \( \Gamma\subset G \) a cocompact torsion-free lattice so that \( X = \Gamma\backslash G/H \) is a compact locally symmetric space of noncompact type. Denote by \(D_X\) the (spin) Dirac operator on \(X\) and by \(\alpha(X)\in KO_*(C^*_r(\Gamma))\) its Rosenberg index class.  

While the existence or nonexistence of positive scalar curvature (PSC) metrics on \(X\) is largely governed by this higher index, very little is explicitly known about how \(\alpha(X)\) depends on the representation-theoretic structure of \(G\).  

**Problem:**  
Find an explicit formula or computable expression for the Rosenberg index class \(\alpha(X)\in KO_*(C^*_r(\Gamma))\) of \(X=\Gamma\backslash G/H\) directly in terms of representation-theoretic data of \(G\); specifically, determine whether the contribution of the discrete series representations (or more generally, \(K\)-types appearing in the Plancherel decomposition of \(L^2(G)\)) can be arranged to give a canonical description of \(\alpha(X)\).  

In particular, investigate the conjectural link:  
\[
\alpha(X)\;\stackrel{?}{=}\;\mathrm{index}_G(D_{G/H})\;\in K_0(C_r^*(G))
\]
under the Baum–Connes assembly map \(K_*(B\Gamma)\to K_*(C_r^*\Gamma)\), and determine if explicit representation-theoretic multiplicities in the \(L^2(G)\) decomposition control the nonvanishing of \(\alpha(X)\).  

Why is it a "good" problem:  
This problem isolates the deep but currently opaque interface between geometric index theory and representation theory. Unlike general PSC nonexistence results (which are established via abstract K-theoretic arguments), this question demands an *explicit* analytic–representation-theoretic realization of the higher index invariant for locally symmetric spaces. The problem is conceptually simple but potentially transformative: resolving it would concretely intertwine the global differential geometry of locally symmetric spaces with the Plancherel formula and the unitary dual of semisimple Lie groups, producing tangible analytic expressions for index-theoretic obstructions. It could reveal exact mechanisms by which representation-theoretic data encode topological curvature obstructions, thereby deepening the link among spin geometry, noncommutative topology, and harmonic analysis on reductive groups.